\subsection{Multi-Domain Systems — Electromechanical, Electrochemical, and Fluidic Systems}

Control systems rarely exist in isolation within a single energy domain. The most interesting and practical systems invariably involve the transfer and conversion of energy across multiple domains—electrical to mechanical, chemical to electrical, hydraulic to mechanical, and countless other combinations. Understanding how to model these multi-domain systems is essential for any control engineer, as it reveals the fundamental unity underlying what might appear to be disparate physical phenomena. The remarkable insight that governs all multi-domain modeling is this: the same mathematical structures and conservation laws apply across domains, even when the physical quantities themselves differ dramatically. A differential equation describing the behavior of an electrical circuit can, with appropriate substitutions, describe the motion of a mechanical system or the flow of fluid through pipes. This correspondence, known as the analogy between domains, provides both practical modeling tools and deep physical insight.

The general approach to multi-domain modeling begins with recognizing that every physical system stores and dissipates energy, and that energy is conserved in all transformations. Whether we are dealing with electrons flowing through a wire, masses moving through space, molecules reacting in a vessel, or fluid moving through pipes, we can identify analogous quantities that play corresponding roles. Electrical voltage becomes mechanical force, electrical current becomes mechanical velocity, electrical resistance becomes mechanical damping, and so forth. These analogies are not mere mathematical curiosities; they reflect deep physical symmetries in how nature stores, dissipates, and transfers energy. By mastering these analogies, the control engineer gains a powerful unified framework applicable to virtually any physical system.

% ------------------------------------------------------------
% Subsection 1.5.5.1: Electromechanical Systems
% ------------------------------------------------------------
.subsubsection{Electromechanical Systems: Motors, Generators, and Transducers}

Electromechanical systems are perhaps the most commonly encountered multi-domain systems in control engineering, appearing in everything from precision positioning stages to industrial robots to automotive subsystems. These systems convert energy between electrical and mechanical forms, and their modeling requires simultaneously applying the governing equations from both domains. The key to successful electromechanical modeling lies in understanding the coupling mechanisms—the physical phenomena that allow energy to flow between domains—and accurately representing these couplings in our mathematical descriptions.

Consider the DC motor, the quintessential electromechanical device, which converts electrical energy into mechanical rotational energy. A DC motor consists of a stator that creates a magnetic field and a rotor (armature) that carries current and experiences forces due to the interaction between this current and the magnetic field. The electrical behavior of the motor is governed by the armature circuit equation, which accounts for the voltage drop across the armature resistance and the induced voltage (back EMF) generated by the rotating armature conductors cutting magnetic flux lines. Simultaneously, the mechanical behavior is governed by Newton's second law applied to rotational motion, which relates the net electromagnetic torque to the moment of inertia and accounts for friction and load torque. These two sets of equations are coupled through two key quantities: the back EMF constant and the torque constant, which are numerically equal in SI units due to energy conservation.

The electrical equation for a DC motor with armature voltage $v_a(t)$, armature current $i_a(t)$, armature resistance $R_a$, and armature inductance $L_a$ is given by the voltage balance equation:
\begin{equation}
v_a(t) = L_a \frac{di_a(t)}{dt} + R_a i_a(t) + e_b(t)
\end{equation}
where $e_b(t)$ is the back EMF, proportional to the angular velocity $\omega_m(t)$:
\begin{equation}
e_b(t) = K_e \omega_m(t)
\end{equation}
The constant $K_e$ is the back EMF constant, typically expressed in volts per radian per second. The mechanical equation, applying Newton's second law to rotation, is:
\begin{equation}
J \frac{d\omega_m(t)}{dt} = \tau_e(t) - \tau_L(t) - B \omega_m(t)
\end{equation}
where $J$ is the total moment of inertia reflected to the motor shaft, $\tau_e(t)$ is the electromagnetic torque developed by the motor, $\tau_L(t)$ is any external load torque, and $B$ is the viscous friction coefficient. The electromagnetic torque is proportional to the armature current:
\begin{equation}
\tau_e(t) = K_t i_a(t)
\end{equation}
where $K_t$ is the torque constant, typically expressed in newton-meters per ampere. The profound result from energy conservation is that $K_e = K_t$ when both are expressed in SI units, a relationship that can be verified by considering the power balance in the motor.

The power delivered to the motor electrically is $P_e = v_a i_a$, while the mechanical power output is $P_m = \tau_e \omega_m$. In steady state with constant velocity, the electrical power input equals the sum of power dissipated in resistance ($i_a^2 R_a$) and mechanical power output. Substituting our relations:
\begin{equation}
P_e = (e_b + i_a R_a) i_a = e_b i_a + i_a^2 R_a = K_e \omega_m i_a + i_a^2 R_a
\end{equation}
\begin{equation}
P_m = \tau_e \omega_m = K_t i_a \omega_m
\end{equation}
For these expressions to be consistent with energy conservation, we must have $K_e = K_t$. This equality is not an accident—it reflects the fundamental symmetry between electrical and mechanical energy conversion, known as the principle of electromechanical reciprocity.

A complete electromechanical model must also account for the gear reduction commonly used between motors and loads. Gears provide mechanical advantage, trading rotational speed for torque, but they also introduce dynamics through the reflected inertia and compliance of the gear teeth. When a motor with torque $\tau_m$ and speed $\omega_m$ drives a load through a gear ratio $N = \omega_m / \omega_L = \tau_L / \tau_m$ (where subscripts $m$ and $L$ denote motor and load quantities respectively), the reflected inertia at the motor shaft becomes $J_{ref} = J_L / N^2$, and the reflected load torque becomes $\tau_{L,ref} = \tau_L / N$. The gear ratio thus has a squared effect on reflected inertia, meaning high gear ratios can dramatically reduce the effective inertia seen by the motor, improving dynamic response at the cost of reduced speed and increased torque demands on the motor.

The complete state-space representation for a DC motor driving an inertial load through a gear reduction incorporates the electrical dynamics (from the inductor), the mechanical dynamics (from the inertia), and the algebraic constraints from the gear ratio. Choosing the state variables as armature current $i_a$ and motor angular velocity $\omega_m$, we obtain:
\begin{equation}
\frac{d}{dt}\begin{bmatrix} i_a \\ \omega_m \end{bmatrix} = \begin{bmatrix} -\frac{R_a}{L_a} & -\frac{K_e}{L_a} \\ \frac{K_t}{J_{eq}} & -\frac{B_{eq}}{J_{eq}} \end{bmatrix} \begin{bmatrix} i_a \\ \omega_m \end{bmatrix} + \begin{bmatrix} \frac{1}{L_a} \\ 0 \end{bmatrix} v_a + \begin{bmatrix} 0 \\ -\frac{1}{J_{eq}} \end{bmatrix} \tau_L
\end{equation}
where $J_{eq} = J_m + J_L / N^2$ and $B_{eq} = B_m + B_L / N^2$ represent the equivalent inertia and friction at the motor shaft. This state-space model provides the foundation for controller design and analysis of DC motor systems.

Generators are the inverse of motors, converting mechanical energy into electrical energy. The modeling approach is identical—we simply interpret the equations differently. When mechanical torque drives the rotor, the back EMF becomes the output voltage, and the current flows out of the machine rather than in. The same mathematical structure describes both devices, which is why many practical machines can operate as either motors or generators depending on the direction of energy flow. This bidirectional capability is essential in applications such as regenerative braking in electric vehicles, where the motor acts as a generator during deceleration, converting kinetic energy back to electrical energy that can be stored in the battery.

\begin{example}{DC Motor Start-Up Transient Analysis}
Consider a DC motor with the following parameters: armature resistance $R_a = 2\ \Omega$, armature inductance $L_a = 0.01\ \text{H}$, back EMF constant $K_e = 0.1\ \text{V·s/rad}$, torque constant $K_t = 0.1\ \text{N·m/A}$, moment of inertia $J = 0.01\ \text{kg·m}^2$, and viscous friction coefficient $B = 0.001\ \text{N·m·s/rad}$. The motor is initially at rest and is suddenly connected to a 24 V source at $t = 0$. We wish to determine the time evolution of the armature current, motor speed, and developed torque.

The system of differential equations governing the motor behavior is:
\begin{equation}
L_a \frac{di_a}{dt} = v_a - R_a i_a - K_e \omega_m
\end{equation}
\begin{equation}
J \frac{d\omega_m}{dt} = K_t i_a - B \omega_m
\end{equation}
Substituting the numerical values:
\begin{equation}
0.01 \frac{di_a}{dt} = 24 - 2i_a - 0.1\omega_m
\end{equation}
\begin{equation}
0.01 \frac{d\omega_m}{dt} = 0.1i_a - 0.001\omega_m
\end{equation}
Simplifying by dividing both sides of each equation by the coefficient of the derivative term:
\begin{equation}
\frac{di_a}{dt} = 2400 - 200i_a - 10\omega_m
\end{equation}
\begin{equation}
\frac{d\omega_m}{dt} = 10i_a - 0.1\omega_m
\end{equation}

To solve this system of first-order linear differential equations, we can express it in matrix form $\dot{\mathbf{x}} = A\mathbf{x} + \mathbf{b}$, where $\mathbf{x} = [i_a, \omega_m]^T$:
\begin{equation}
\dot{\mathbf{x}} = \begin{bmatrix} -200 & -10 \\ 10 & -0.1 \end{bmatrix} \mathbf{x} + \begin{bmatrix} 2400 \\ 0 \end{bmatrix}
\end{equation}
The equilibrium point occurs when $\dot{\mathbf{x}} = 0$, which gives:
\begin{equation}
\begin{bmatrix} -200 & -10 \\ 10 & -0.1 \end{bmatrix} \begin{bmatrix} i_{a,ss} \\ \omega_{ss} \end{bmatrix} + \begin{bmatrix} 2400 \\ 0 \end{bmatrix} = 0
\end{equation}
Solving the second row first (which is simpler): $10i_{a,ss} - 0.1\omega_{ss} = 0$, so $\omega_{ss} = 100i_{a,ss}$. Substituting into the first row: $-200i_{a,ss} - 10(100i_{a,ss}) + 2400 = 0$, which simplifies to $-1200i_{a,ss} + 2400 = 0$, giving $i_{a,ss} = 2\ \text{A}$ and $\omega_{ss} = 200\ \text{rad/s}$.

To find the transient response, we determine the eigenvalues of the system matrix $A$. The characteristic equation is:
\begin{equation}
\det(A - \lambda I) = \det\begin{bmatrix} -200 - \lambda & -10 \\ 10 & -0.1 - \lambda \end{bmatrix} = (\lambda + 200)(\lambda + 0.1) + 100 = 0
\end{equation}
Expanding: $\lambda^2 + 200.1\lambda + 20 + 100 = \lambda^2 + 200.1\lambda + 120 = 0$. Using the quadratic formula:
\begin{equation}
\lambda = \frac{-200.1 \pm \sqrt{200.1^2 - 4 \cdot 1 \cdot 120}}{2} = \frac{-200.1 \pm \sqrt{40040.01 - 480}}{2} = \frac{-200.1 \pm \sqrt{39560.01}}{2}
\end{equation}
\begin{equation}
\lambda = \frac{-200.1 \pm 198.90}{2}
\end{equation}
The two eigenvalues are $\lambda_1 = -199.5\ \text{s}^{-1}$ and $\lambda_2 = -0.6\ \text{s}^{-1}$. Both eigenvalues are real and negative, indicating a stable, overdamped response with two distinct time constants: $\tau_1 = 1/199.5 \approx 0.005\ \text{s}$ (the electrical time constant dominated by inductance) and $\tau_2 = 1/0.6 \approx 1.67\ \text{s}$ (the mechanical time constant dominated by inertia).

The general solution for each state variable is the sum of the particular (steady-state) solution and the homogeneous (transient) solution:
\begin{equation}
i_a(t) = i_{a,ss} + C_1 e^{\lambda_1 t} + C_2 e^{\lambda_2 t} = 2 + C_1 e^{-199.5t} + C_2 e^{-0.6t}
\end{equation}
\begin{equation}
\omega_m(t) = \omega_{ss} + D_1 e^{\lambda_1 t} + D_2 e^{\lambda_2 t} = 200 + D_1 e^{-199.5t} + D_2 e^{-0.6t}
\end{equation}
Applying the initial conditions $i_a(0) = 0$ and $\omega_m(0) = 0$:
\begin{equation}
i_a(0) = 2 + C_1 + C_2 = 0 \quad\Rightarrow\quad C_1 + C_2 = -2
\end{equation}
\begin{equation}
\omega_m(0) = 200 + D_1 + D_2 = 0 \quad\Rightarrow\quad D_1 + D_2 = -200
\end{equation}
Using the derivative relationships from the differential equations evaluated at $t = 0$:
\begin{equation}
\left.\frac{di_a}{dt}\right|_{t=0} = 2400 - 200i_a(0) - 10\omega_m(0) = 2400\ \text{A/s}
\end{equation}
\begin{equation}
\left.\frac{d\omega_m}{dt}\right|_{t=0} = 10i_a(0) - 0.1\omega_m(0) = 0\ \text{rad/s}^2
\end{equation}
The derivative of the current solution gives:
\begin{equation}
\left.\frac{di_a}{dt}\right|_{t=0} = -199.5C_1 - 0.6C_2 = 2400
\end{equation}
The derivative of the speed solution gives:
\begin{equation}
\left.\frac{d\omega_m}{dt}\right|_{t=0} = -199.5D_1 - 0.6D_2 = 0
\end{equation}
Solving these four equations simultaneously (the initial conditions and derivatives) yields $C_1 \approx -1.994$, $C_2 \approx -0.006$, $D_1 \approx 0.995$, and $D_2 \approx -200.995$. The complete solution is:
\begin{equation}
i_a(t) = 2 - 1.994e^{-199.5t} - 0.006e^{-0.6t}\ \text{A}
\end{equation}
\begin{equation}
\omega_m(t) = 200 + 0.996e^{-199.5t} - 201.0e^{-0.6t}\ \text{rad/s}
\end{equation}
The developed torque is $\tau_e(t) = K_t i_a(t) = 0.1i_a(t)\ \text{N·m}$. This analysis reveals that the electrical transient (with time constant 5 ms) is much faster than the mechanical transient (with time constant 1.67 s). The armature current rises rapidly to nearly its final value, then the motor speed gradually increases as the inertia is overcome against friction. The initial current surge is limited by the back EMF, which builds up as the motor accelerates. The developed torque peaks early (at approximately 0.24 N·m) when current is maximum and then decreases as speed increases and current approaches its steady-state value.
\end{example}

Transducers are devices that convert energy between domains in a way that enables measurement of physical quantities. Strain gauges convert mechanical deformation to electrical resistance changes, thermocouples convert temperature differences to voltage differences, piezoelectric sensors convert mechanical stress to electrical charge, and accelerometers convert acceleration to electrical signals. Each transducer has its own characteristic modeling considerations, but all share the common feature of coupling two domains. A piezoelectric accelerometer, for instance, exhibits both mechanical behavior (the seismic mass's response to acceleration) and electrical behavior (the charge generated by the piezoelectric material). The equivalent circuit model of such a device typically includes a mechanical mass-spring system whose output is modeled as a current source proportional to acceleration, connected in parallel with the capacitance of the piezoelectric element. Understanding the complete coupled model is essential for proper signal conditioning and accurate measurement.

% ------------------------------------------------------------
% Subsection 1.5.5.2: Fluidic and Hydraulic Systems
% ------------------------------------------------------------
.subsubsection{Hydraulic and Fluidic Systems: The Flow-Pressure Analogy}

Hydraulic systems present a particularly elegant example of multi-domain modeling because their behavior maps directly onto electrical circuit theory through well-defined analogies. This correspondence is not merely mathematical convenience—it reflects the deep similarity in how fluids and electricity store, dissipate, and transport energy. Two primary analogies exist: the force-voltage (impedance) analogy and the force-current (mobility) analogy. The force-voltage analogy, which we will develop here, maps fluid pressure to electrical voltage and volumetric flow rate to electrical current. This mapping has the intuitive appeal of treating pressure as the "driving force" for flow, analogous to voltage driving current through a circuit.

In the force-voltage analogy, the fundamental quantities are defined as follows. Pressure $p(t)$, measured in pascals (Pa), corresponds to voltage $v(t)$, measured in volts (V). Volumetric flow rate $Q(t)$, measured in cubic meters per second (m³/s), corresponds to current $i(t)$, measured in amperes (A). This mapping is physically grounded in the observation that both pressure difference and voltage difference serve as the driving potential for flow of their respective conserved quantities (fluid volume and electrical charge). The conservation laws for both domains become statements of Kirchhoff's laws: the net flow into any junction must be zero (Kirchhoff's Current Law), and the sum of potential differences around any closed loop must be zero (Kirchhoff's Voltage Law).

Fluid resistance arises from viscous dissipation as fluid flows through pipes, valves, or orifices. For laminar flow through a cylindrical pipe, the resistance is given by the Hagen-Poiseuille equation, which relates the pressure drop to the flow rate through the pipe's geometry and fluid viscosity. For turbulent flow or flow through restrictions, the resistance relationship becomes nonlinear, with pressure drop proportional to the square of flow rate. In electrical terms, a linear fluid resistance corresponds to an electrical resistor obeying Ohm's law, while a nonlinear restriction corresponds to a nonlinear resistor. The power dissipated in fluid resistance is $P = p Q$, analogous to electrical power $P = v i$, and this power represents irreversible conversion to heat in both cases.

The mathematical relationship for laminar flow through a circular pipe of length $L$ and radius $r$ is:
\begin{equation}
Q = \frac{\pi r^4}{8\mu L} \Delta p
\end{equation}
where $\mu$ is the dynamic viscosity of the fluid. The resistance $R_f = \Delta p / Q$ is therefore:
\begin{equation}
R_f = \frac{8\mu L}{\pi r^4}
\end{equation}
This relationship reveals the dramatic sensitivity of fluid resistance to pipe radius—a halving of the radius increases resistance by a factor of sixteen. This nonlinear scaling explains why small restrictions in hydraulic systems can create significant pressure drops even at modest flow rates.

Fluid capacitance represents the ability of a volume to store fluid under pressure, analogous to a capacitor storing electrical charge under voltage. A simple example is a flexible bladder or accumulator, where fluid entering the volume compresses a gas, storing energy in the compression. For a gas-charged accumulator following isothermal compression, the capacitance $C_f$ relates the change in stored volume to the change in pressure:
\begin{equation}
C_f = \frac{dV}{dp} = \frac{V_0}{p_0}
\end{equation}
where $V_0$ is the initial gas volume and $p_0$ is the initial gas pressure (absolute). For rigid-walled tanks with free surfaces (such as an open reservoir), the capacitance is simply the cross-sectional area $A$ of the tank surface, since a change in fluid level $h$ corresponds to a volume change $dV = A dh$ and a pressure change $dp = \rho g dh$ at depth, where $\rho$ is the fluid density and $g$ is gravitational acceleration. In this case, the relationship between flow rate and pressure is:
\begin{equation}
Q = C_f \frac{dp}{dt}
\end{equation}
where $C_f = A / (\rho g)$ is the effective capacitance of the open tank. This is directly analogous to the capacitor equation $i = C dv/dt$.

Fluid inductance represents the inertial properties of fluid mass in motion, analogous to electrical inductance representing the inertia of moving charges. When the flow rate through a conduit changes, the fluid mass must be accelerated or decelerated, requiring a pressure difference according to Newton's second law. For a fluid of density $\rho$ moving through a pipe of cross-sectional area $A$ and length $L$, the inductance $L_f$ is:
\begin{equation}
L_f = \frac{\rho L}{A}
\end{equation}
The governing equation relating pressure drop to flow rate is:
\begin{equation}
p = L_f \frac{dQ}{dt}
\end{equation}
This equation states that a changing flow rate requires a proportional pressure difference to accelerate the fluid mass. In electrical terms, this is identical to the voltage across an inductor: $v = L di/dt$.

The complete analogy between electrical and hydraulic quantities is summarized in Table~\ref{tab:domain_analogy}, which provides a comprehensive mapping between electrical and hydraulic systems. This table serves as a reference for translating electrical circuit concepts into hydraulic contexts and vice versa. Understanding these equivalences allows control engineers to apply circuit analysis techniques—nodal analysis, mesh analysis, transfer function methods, frequency response analysis—to hydraulic systems with confidence that the mathematical results will accurately predict physical behavior.

\begin{table}[h]
\centering
\bgroup
\def\arraystretch{1.2}
\begin{tabular}{|l|l|l|}
\hline
\textbf{Quantity} & \textbf{Electrical Domain} & \textbf{Hydraulic Domain} \\ \hline
Potential (across variable) & Voltage $v$ [V] & Pressure $p$ [Pa] \\ \hline
Flow (through variable) & Current $i$ [A] & Volumetric flow $Q$ [m³/s] \\ \hline
Dissipative element & Resistance $R$ [$\Omega$] & Fluid resistance $R_f$ [Pa·s/m³] \\ \hline
Storage element (energy) & Capacitance $C$ [F] & Fluid capacitance $C_f$ [m³/Pa] \\ \hline
Storage element (momentum) & Inductance $L$ [H] & Fluid inductance $L_f$ [Pa·s²/m³] \\ \hline
Energy stored in capacitance & $E_C = \frac{1}{2}CV^2$ & $E_C = \frac{1}{2}C_f p^2$ \\ \hline
Energy stored in inductance & $E_L = \frac{1}{2}LI^2$ & $E_L = \frac{1}{2}L_f Q^2$ \\ \hline
Power dissipation & $P = vi = i^2R$ & $P = pQ = Q^2 R_f$ \\ \hline
Kirchhoff's Current Law & $\sum i = 0$ at node & $\sum Q = 0$ at junction \\ \hline
Kirchhoff's Voltage Law & $\sum v = 0$ around loop & $\sum p = 0$ around loop \\ \hline
Impedance & $Z = V/I$ & $Z_f = p/Q$ \\ \hline
\end{tabular}
\egroup
\\[0.5em]
\caption{Analogous quantities between electrical and hydraulic domains using the force-voltage analogy. This mapping enables circuit analysis techniques to be applied to hydraulic systems.\label{tab:domain_analogy}}
\end{table}

The practical importance of hydraulic systems in control applications cannot be overstated. Hydraulic actuators provide high power density—much higher than electric motors of comparable size—making them essential in heavy machinery, aerospace systems, and industrial presses. The modeling approach we have developed allows the control engineer to predict the dynamic response of hydraulic systems, design appropriate controllers, and understand the limitations imposed by fluid properties and system geometry. A hydraulic servo system, for instance, can be analyzed using the same frequency response techniques applied to electrical amplifiers, with the understanding that the hydraulic components introduce specific dynamics related to fluid compressibility, line capacitance, and inertial effects.

\begin{example}{Hydraulic Circuit Transient Analysis}
Consider a simple hydraulic circuit consisting of a pump maintaining constant pressure, connected through a pipe to an accumulator (a gas-charged pressure vessel), and then through a valve to a hydraulic cylinder. We wish to determine how the flow rate and pressure at the cylinder respond when the valve is suddenly opened, allowing fluid to enter the cylinder and extend a piston against a spring load. The system parameters are: supply pressure $p_s = 10^6\ \text{Pa}$, pipe resistance $R_p = 10^8\ \text{Pa·s/m}^3$, pipe inductance $L_p = 10^5\ \text{Pa·s}^2/\text{m}^3$, accumulator capacitance $C_a = 10^{-8}\ \text{m}^3/\text{Pa}$, valve resistance $R_v = 2 \times 10^8\ \text{Pa·s/m}^3$, cylinder area $A_c = 10^{-3}\ \text{m}^2$, piston mass $m = 100\ \text{kg}$, and spring constant $k = 10^5\ \text{N/m}$.

The system can be modeled as a hydraulic network with three primary elements: the supply network (pipe and accumulator), the valve, and the load (cylinder with spring). The governing equations begin with the accumulator, which stores fluid under pressure and thus behaves as a hydraulic capacitor:
\begin{equation}
Q_p - Q_v = C_a \frac{dp_a}{dt}
\end{equation}
where $Q_p$ is the flow from the pump, $Q_v$ is the flow through the valve, and $p_a$ is the accumulator pressure. The pipe between the accumulator and valve introduces both resistance and inductance:
\begin{equation}
p_a - p_v = R_p Q_p + L_p \frac{dQ_p}{dt}
\end{equation}
where $p_v$ is the pressure at the valve inlet. The valve has purely resistive behavior (assuming it is fully opened and operating in a linear regime):
\begin{equation}
p_v - p_c = R_v Q_v
\end{equation}
where $p_c$ is the pressure at the cylinder inlet. The cylinder converts fluid flow to mechanical motion, with the flow rate related to piston velocity by $Q_v = A_c \dot{x}$, where $x$ is the piston position. The mechanical dynamics of the piston are governed by Newton's second law:
\begin{equation}
A_c p_c = m \ddot{x} + kx + B \dot{x}
\end{equation}
where we include viscous damping $B$ with value $1000\ \text{N·s/m}$.

For simplicity, we assume the pump can maintain constant flow $Q_p = 10^{-4}\ \text{m}^3/\text{s}$ and that the system is initially at steady state with the valve closed ($Q_v = 0$). At $t = 0$, the valve opens, and we analyze the transient response. Under steady-state conditions with the valve closed, the accumulator pressure equals the supply pressure: $p_a = p_s = 10^6\ \text{Pa}$. With $Q_p = 10^{-4}\ \text{m}^3/\text{s}$, the pressure drop across the pipe is $R_p Q_p = 10^8 \times 10^{-4} = 10^4\ \text{Pa}$, so $p_v = p_s - R_p Q_p = 990,000\ \text{Pa}$.

After the valve opens, the system reaches a new steady state where $Q_p = Q_v = 10^{-4}\ \text{m}^3/\text{s}$ (since the accumulator capacitance maintains constant pressure). The pressure drop across the valve is $R_v Q_v = 2 \times 10^8 \times 10^{-4} = 20,000\ \text{Pa}$. Thus, $p_c = p_v - R_v Q_v = 990,000 - 20,000 = 970,000\ \text{Pa}$. The steady-state piston force from pressure is $A_c p_c = 10^{-3} \times 970,000 = 970\ \text{N}$. This force must balance the spring force at equilibrium displacement: $kx_{ss} = 970\ \text{N}$, giving $x_{ss} = 970 / 10^5 = 0.0097\ \text{m} = 9.7\ \text{mm}$.

To find the transient response, we linearize the system about the operating point and derive the transfer function relating piston position to the valve opening. The combined dynamics are dominated by two time constants: the hydraulic time constant associated with the accumulator capacitance and pipe inductance, and the mechanical time constant associated with the piston mass and damping. The hydraulic time constant is approximately $\tau_h = \sqrt{L_p C_a} = \sqrt{10^5 \times 10^{-8}} = \sqrt{10^{-3}} = 0.032\ \text{s}$. The mechanical time constant is $\tau_m = m/B = 100/1000 = 0.1\ \text{s}$. Since the mechanical time constant is larger, the overall system response is dominated by the mechanical dynamics, with the hydraulic subsystem providing a faster response that can be approximated as instantaneous for preliminary analysis.

The complete transfer function relating piston position $X(s)$ to the valve flow coefficient (which changes when the valve opens) is:
\begin{equation}
\frac{X(s)}{K_{valve}} = \frac{A_c / (k s)}{(\tau_m^2 s^2 + 2\zeta_m \tau_m s + 1)(\tau_h^2 s^2 + 2\zeta_h \tau_h s + 1)}
\end{equation}
where $\zeta_m \approx 0.05$ (lightly damped mechanical system) and $\zeta_h \approx 0.5$ (moderately damped hydraulic system). This transfer function predicts an oscillatory transient response with natural frequencies $\omega_{n,m} = 1/\tau_m = 10\ \text{rad/s}$ in the mechanical domain and $\omega_{n,h} = 1/\tau_h = 31.6\ \text{rad/s}$ in the hydraulic domain. The lower mechanical natural frequency dominates the overall response, and the system will exhibit oscillations at approximately $10\ \text{rad/s}$ (about 1.6 Hz) during the transient as the piston and spring exchange energy, gradually settling to the new equilibrium position as the mechanical damping dissipates energy.
\end{example}

% ------------------------------------------------------------
% Subsection 1.5.5.3: Energy Balance and System Integration
% ------------------------------------------------------------
.subsubsection{Energy Balance and System Integration Across Domains}

The concept of energy balance provides a unifying framework for understanding multi-domain systems, ensuring that our models respect the fundamental conservation of energy. In any physical system, the rate at which energy enters through the boundaries must equal the rate at which energy is stored within the system plus the rate at which energy is dissipated. This principle, known as the power balance equation, applies identically across all energy domains and provides a powerful check on the correctness of multi-domain models. When we couple electrical, mechanical, hydraulic, and thermal subsystems into an integrated model, energy balance analysis helps identify errors in the coupling equations and ensures that our controller designs do not violate physical constraints.

For a general multi-domain system, the power balance equation can be written as:
\begin{equation}
\sum_{in} P_{in} = \frac{d}{dt}(E_{stored}) + \sum_{diss} P_{diss}
\end{equation}
where the left side represents the total power input from all sources, the first term on the right represents the rate of change of stored energy (in all its forms—electrical, mechanical, thermal, chemical), and the second term represents the total power dissipated (converted to heat and irreversibly lost). In a well-designed system, all power inputs eventually become either stored energy (temporarily) or dissipated energy (permanently). This accounting must balance at every instant in time.

Consider an electromechanical system consisting of a DC motor driving a pump that moves fluid through a pipe with a restriction. Energy enters as electrical power $P_e = v i$ from the power supply. Some of this power is stored in the magnetic field of the motor's inductance ($d/dt(\frac{1}{2}Li^2)$) and some is dissipated as heat in the armature resistance ($i^2R_a$). The remaining electrical power is converted to mechanical power at the shaft: $P_m = \tau \omega$. The mechanical power is partially stored in the rotational kinetic energy of the rotor ($d/dt(\frac{1}{2}J\omega^2)$) and partially dissipated in friction ($B\omega^2$). The remainder is delivered to the pump, which converts it to hydraulic power $P_h = pQ$. The hydraulic power is partially stored in the kinetic energy of the moving fluid ($d/dt(\frac{1}{2}\rho L Q^2/A)$ for pipe inductance effects) and partially dissipated in viscous friction along the pipe walls and through the restriction ($Q^2 R_f$). At every point in this energy conversion chain, the power balance holds: the power leaving each subsystem equals the power entering the next, accounting for local storage and dissipation.

The integration of multi-domain systems for control purposes requires careful attention to the interfaces between domains. These interfaces—the transducers, actuators, and sensors that connect different domains—are where modeling errors most commonly occur. A motor model that neglects the back EMF, for instance, will incorrectly predict that the current continues to increase indefinitely as the motor accelerates, violating energy conservation. A hydraulic model that neglects fluid compressibility will predict instantaneous pressure changes that are physically impossible, since pressure waves propagate at the speed of sound through the fluid. A thermal model that neglects heat capacity will predict infinitely fast temperature changes. These interface errors often manifest as non-physical behaviors such as negative time constants, unbounded growth, or responses that violate causality.

The cascade structure of most multi-domain systems—the electrical controller driving an electrical amplifier, which drives an electromechanical actuator, which drives a hydraulic mechanism, which ultimately acts on a mechanical load—suggests a hierarchical approach to modeling and control design. Each stage in the cascade can be modeled independently, verified against energy balance, and then connected to the next stage. The bandwidth of each stage typically decreases as we move from electrical to mechanical domains, due to the inherent time constants of each domain. Electrical circuits can respond at megahertz frequencies, electromechanical systems at kilohertz frequencies, hydraulic systems at hundreds of hertz, and thermal systems at millihertz. This cascade of bandwidth limitations has important implications for control design: a controller designed for a thermal system cannot compensate for dynamics occurring at hydraulic or electrical frequencies, and similarly for each domain.

The principles developed in this section extend naturally to other energy domains not explicitly covered here. Chemical systems can be modeled using analogous quantities where chemical potential plays the role of voltage and molar flow rate plays the role of current. Thermal systems use temperature as the potential and heat flow as the flow quantity. Even thermodynamic systems with their inherent irreversibilities can be modeled using network analogies, with thermal resistances representing temperature gradients and thermal capacitances representing heat capacities. The control engineer who masters these analogies gains a powerful toolkit applicable to virtually any physical system, unified by the fundamental laws of physics governing energy storage, dissipation, and transfer.

In summary, multi-domain system modeling requires the control engineer to think across traditional disciplinary boundaries, recognizing the mathematical and physical unity underlying seemingly disparate phenomena. The key steps are: identify the relevant energy domains in the system, select appropriate analogies between domains, write the governing equations for each subsystem using the domain-specific physics, couple the subsystems through appropriate interface equations (verifying energy conservation at each interface), and assemble the complete model using the conservation laws that apply at each level of integration. This systematic approach, grounded in first principles and verified by energy balance, produces models that are physically meaningful, mathematically tractable, and useful for control design.
